% Part: normal-modal-logic
% Chapter: axioms-systems
% Section: logics-proofs

\documentclass[../../../include/open-logic-section]{subfiles}

\begin{document}

\olfileid{nml}{prf}{prf}

\olsection{\usetoken{P}{derivation} અને મોડલ તંત્રો}

પહેલાં આપણે સામાન્ય મોડલ તર્કો માટે !!a{derivation} શું છે તે
વ્યાખ્યાયિત કરીએ. અંદાજે, !!a{derivation} એ !!{formula}sની એવી શ્રેણી છે
જેમાં દરેક !!{element} કાં તો કેટલાંક \emph{સ્વયંસિદ્ધિઓ}માંનું એક
(અથવા તેનો પ્રતિસ્થાપન દૃષ્ટાંત) હોય છે, અથવા અગાઉનાં !!{element}sમાંથી
થોડા અનુમાન-નિયમોમાંના કોઈ નિયમ વડે ફલિત થાય છે. સામાન્ય મોડલ તર્કો
માટે પુનરુક્તિઓના બધા દૃષ્ટાંતો\iftag{prvDiamond}{, \Ax{K} અને~\Dual{}}{ અને~\Ax{K}}
સ્વયંસિદ્ધિઓ ગણાય છે. તેનાથી સૌથી નાનો સામાન્ય મોડલ તર્ક, મોડલ
તંત્ર~$\Log{K}$, મળે છે. બીજા તંત્રો મેળવવા આપણે વધારાનાં
સ્વયંસિદ્ધિઓ ઉમેરી શકીએ. નિયમો હંમેશા મોડસ પોનેન્સ~\MP{} અને
આવશ્યકીકરણ~\Nec છે.

\begin{defn}
  મોડલ તંત્ર $\Log{K} !A_1 \dots !A_n$ અને !!a{formula}~$!B$ આપેલાં
  હોય, ત્યારે $!B$ એ $\Log{K} !A_1 \dots
  !A_n$માં \emph{!!{derivable}} છે, જે
  $\Log{K} !A_1 \dots !A_n \Proves !B$ લખાય છે, જો અને માત્ર જો
  એવાં !!{formula}s $!C_1$, \dots, $!C_k$ હોય કે $!C_k = !B$ અને દરેક
  $!C_i$ કાં તો પુનરુક્તિનો દૃષ્ટાંત હોય, અથવા
  $\Ax{K}$,\iftag{prvDiamond}{ $\Dual$,}{} $!A_1$,
  \dots, $!A_n$ પૈકી કોઈ એકનો દૃષ્ટાંત હોય, અથવા અગાઉનાં
  !!{formula}sમાંથી \MP{} કે~\Nec નિયમ વડે ફલિત થાય.
\end{defn}

નીચેનો પ્રસ્તાવ~$!B$ની $\Sigma$-!!{derivation} દર્શાવીને
$!B \in \Sigma$ સાબિત કરવા દે છે.

\begin{prop}
  $\Log{K} !A_1 \dots !A_n = \Setabs{!B}{\Log{K} !A_1 \dots !A_n \Proves !B}$.
\end{prop}

\begin{proof}
  !!{derivation}sની લંબાઈ પર આગમન વાપરીને આપણે બતાવીએ છીએ કે
  $\Setabs{!B}{\Log{K} !A_1 \dots !A_n \Proves !B} \subseteq
  \Log{K} !A_1 \dots !A_n$.

  જો~$!B$ની !!{derivation}ની લંબાઈ~$1$ હોય, તો તેમાં એક જ !!{formula}
  છે. તે !!{formula} અગાઉનાં સૂત્રોમાંથી \MP{} કે \Nec વડે ફલિત થઈ
  શકે નહિ; એટલે તે પુનરુક્તિનો દૃષ્ટાંત, \Ax{K}નો દૃષ્ટાંત,\iftag{prvDiamond}{
  $\Dual$નો દૃષ્ટાંત,}{} અથવા $!A_1$, \dots,~$!A_n$ પૈકી કોઈ એકનો
  દૃષ્ટાંત હોવું જોઈએ. પણ $\Log{K}!A_1\dots!A_n$માં આ બધાં હોય છે;
  તેથી $!B \in \Log{K}!A_1 \dots !A_n$.

  જો $!B$ની !!{derivation}ની લંબાઈ~$> 1$ હોય, તો $!B$ એ !!{derivation}ની
  છેલ્લી પંક્તિમાં ન આવતાં અગાઉનાં !!{formula}sમાંથી \MP{} કે \Nec{}
  વડે પણ મેળવાયેલું હોઈ શકે છે. જો $!B$ એ $!C$ અને $!C \lif !B$માંથી
  (\MP વડે) ફલિત થતું હોય, તો આગમન-પરિકલ્પના અનુસાર $!C$ અને $!C \lif !B \in
  \Log{K}!A_1 \dots !A_n$. પણ દરેક મોડલ તર્ક મોડસ પોનેન્સ હેઠળ બંધ
  છે; તેથી $!B \in \Log{K}!A_1 \dots
  !A_n$. જો $!B \equiv \Box !C$ એ $!C$માંથી \Nec વડે ફલિત થતું હોય,
  તો આગમન-પરિકલ્પના અનુસાર $!C
  \in \Log{K}!A_1 \dots !A_n$. પણ દરેક સામાન્ય મોડલ તર્ક $\Nec$
  હેઠળ બંધ છે; તેથી $!B \in
  \Log{K}!A_1\dots!A_n$.

  વિપરીત સમાવેશ સાબિત કરવા આપણે બતાવીએ છીએ કે
  $\Sigma = \Setabs{!B}{\Log{K} !A_1 \dots !A_n \Proves !B }$ એ
  $!A_1$, \dots, $!A_n$ના બધા દૃષ્ટાંતો ધરાવતો સામાન્ય મોડલ તર્ક છે.
  પછી વ્યાખ્યા મુજબ $\Log{K} !A_1 \dots !A_n$ સૌથી નાનો એવો તર્ક
  છે, એ અવલોકન વાપરીએ છીએ.
  \begin{enumerate}
    \item દરેક પુનરુક્તિ~$!B$ પુનરુક્તિનો દૃષ્ટાંત હોવાથી
      $\Log{K}!A_1\dots!A_n \Proves !B$; તેથી $\Sigma$ બધી
      પુનરુક્તિઓ ધરાવે છે.
    \item જો $\Log{K}!A_1\dots!A_n \Proves !C$ અને
      $\Log{K}!A_1\dots!A_n \Proves !C \lif !B$ હોય, તો
      $\Log{K}!A_1\dots!A_n \Proves !B$: $!C$ની !!{derivation}ને
      $!C \lif !B$ની નિષ્પત્તિ સાથે જોડો અને અંતે~$!B$ની પંક્તિ
      ઉમેરો. છેલ્લી પંક્તિ \MP{}થી સમર્થિત છે. તેથી $\Sigma$ મોડસ
      પોનેન્સ હેઠળ બંધ છે.
    \item જો $!B$ની !!a{derivation} હોય, તો $!B$ના દરેક પ્રતિસ્થાપન
      દૃષ્ટાંતની પણ !!a{derivation} હોય છે: !!{derivation}માંના દરેક
      !!{formula} પર એ પ્રતિસ્થાપન લાગુ કરો. (અભ્યાસપ્રશ્ન: !!{derivation}sની
      લંબાઈ પર આગમનથી સાબિત કરો કે પરિણામ પણ યોગ્ય !!{derivation} છે.)
      તેથી $\Sigma$ એકસાથે પ્રતિસ્થાપન હેઠળ બંધ છે. (હવે આપણે
      સ્થાપ્યું છે કે $\Sigma$ મોડલ તર્કની બધી શરતો પૂરી કરે છે.)
    \item આપણી પાસે $\Log{K}!A_1\dots!A_n \Proves \Ax{K}$ છે;
      તેથી $\Ax{K} \in \Sigma$.\sourcecorrection{OLMD-024}{મૂળ નિષ્કર્ષમાં $K \in \Sigma$ છે, પરંતુ અહીં \Sigma સૂત્રોનો ગણ છે અને પૂર્વપક્ષમાં \Ax{K} છે; નિષ્કર્ષમાં પણ એ જ સૂત્રચિહ્ન મૂક્યું છે.}
    \tagitem{prvDiamond}{આપણી પાસે
      $\Log{K}!A_1\dots!A_n \Proves \Dual$ છે; તેથી $\Dual \in \Sigma$.}{}
    \item જો $\Log{K}!A_1\dots!A_n \Proves !C$ હોય, તો વધારાની
      પંક્તિ~$\Box !C$ને~\Nec વડે સમર્થન મળે છે. તેથી $\Sigma$
      એ~\Nec હેઠળ બંધ છે; એટલે $\Sigma$ સામાન્ય છે.
  \end{enumerate}
\end{proof}

\end{document}
